% Shortened Chapter 3, part 9 of 10: section 3.9.
% Include after section 3.8; do not include alongside the original section 3.9.
% Subsection order is load-bearing: chapter2.tex cites Section 3.9.5 (metrics)
% and Section 3.9.7 (implications) by number. All eight slots are retained.

\section{Training Under Severe Class Imbalance}
\label{sec:imbalance-review}

This section covers the training configuration held constant across all twelve ablated configurations, and which \autoref{sec:corpus-implications} identified as forced by the corpus rather than chosen freely: the handling of a 4 : 1.3 : 1 class imbalance.

\textbf{A deliberate exclusion.} Two topics the project's reading list groups here --- Grad-CAM++ visual auditing and adversarial identity disentanglement --- are \textbf{not reviewed}, because neither is part of the implemented system. The loss module records that ``the Cross-Batch Memory (XBM) queue, and the adversarial identity loss --- are intentionally removed'', and no gradient-based visualisation exists anywhere in the codebase. Both are declared as proposed further work.

\subsection{The imbalance is a property of the phenomenon, not of the sampling}
\label{sec:imbalance-property}

\autoref{sec:label-taxonomy} and \autoref{sec:original-baseline} established the class counts and the always-Negative floor; what matters here is why the skew exists. It is not an artefact of careless collection. Yan et al.~\cite{yan2014} are explicit that it follows from what can be elicited in a laboratory: ``some types of facial expressions are difficult to elicit in laboratory situations, thus the samples in different categories distributed unequally, e.g., there are 60 disgust samples but only 7 sadness samples.'' Nor is it specific to CASME II --- the MEGC composite shows the same pattern (\autoref{sec:megc-metrics}). Skew must be coped with; it cannot be collected away.

\subsection{Three places to intervene}
\label{sec:imbalance-three-places}

The reviewed corpus corrects for imbalance at three distinct points, worth separating because this thesis intervenes at all three and the interaction between two of them turned out to matter.

\begin{enumerate}
	\item \textbf{At the data}, by changing how often each class is presented during training --- resampling or augmentation.
	\item \textbf{At the loss}, by changing how much each example contributes to the gradient --- class weighting or focal re-weighting.
	\item \textbf{At the metric}, by changing how performance is scored --- macro-averaged rather than sample-averaged measures.
\end{enumerate}

\subsection{Loss-level correction: focal loss}
\label{sec:imbalance-focal}

Zhao et al.~\cite{zhaoS2021} bring focal loss into micro-expression recognition ``to alleviate the inefficient model training caused by the class imbalance in the MEs Datasets'', since cross-entropy ``could be biased toward particular emotions that constitute a larger portion of the training set''. They adopt the formulation developed for one-stage object detection, where the imbalance is between foreground and background regions:

\[
	\text{FL}(p_t) = -\alpha_t (1 - p_t)^{\gamma}\log p_t
\]

where $p_t$ is the model's probability for the true class, $\gamma$ ``the balance factor for loss'' and $\alpha_t$ ``the weight balance factor for samples'' (\autoref{sec:focal-loss}); under skew, the examples it up-weights are disproportionately the minority classes. Adapting it to multi-class classification, Zhao et al. report that ``$\gamma$ is set as 2 in practice, and $\alpha$ is treated as a hyper-parameter to set by cross validation.''

The division of labour there matters: $\gamma$ re-weights by \textit{difficulty} and is treated as a constant, $\alpha$ by \textit{class frequency} and is treated as dataset-specific. Conflating them is where this project went wrong (\autoref{sec:imbalance-double-correction}).

\subsection{Data-level correction: resampling and augmentation}
\label{sec:imbalance-balanced-loss}

Xia et al.~\cite{xia2020a} address the same problem at the data end and pair the two interventions, using ``temporal data augmentation strategies as well as a balanced loss \ldots jointly'' to overcome ``limited and imbalanced training samples''. The pairing is not left unexamined: their subsection \textit{``The Impact of Data Augmentation and Balanced Loss''} removes each correction in turn from the joint system and reports the balanced loss improving performance ``slightly'', with one leave-one-subject-out cell on SMIC where removing it scores higher. This is the corpus's only controlled look at composing two corrections, and \autoref{sec:imbalance-double-correction} records a far larger effect in the same direction.

\subsection{Metric-level correction}
\label{sec:imbalance-metrics}

\autoref{sec:megc-metrics} defines UF1 and the case for macro-averaged metrics. Liong et al.~\cite{liong2016} add the argument from the subject side rather than the class side, reporting F1, precision and recall because these ``provide a more meaningful perspective than accuracy rates when the datasets used are naturally imbalanced since each subject has a different number of video samples.'' The two arguments compound on CASME II, where classes are skewed \textbf{and} per-subject clip counts range from 1 to 33 (\autoref{tab:corpus-threats}c).

\subsection{Limitations}
\label{sec:imbalance-limitations}

\textbf{The interaction between corrections is examined once, and only at small magnitude.} Xia et al.~\cite{xia2020a} ablate augmentation and balanced loss out of the joint system and report the balanced loss helping only ``slightly'' --- at most about two points across their twelve cells, and in one cell costing accuracy; Zhao et al.~\cite{zhaoS2021} apply focal loss alongside a large per-sample augmentation, but without class-proportional resampling and without ablating the pairing. What no reviewed paper reports is the pairing failing outright.

\textbf{$\gamma$ is inherited rather than tuned.} Zhao et al.~\cite{zhaoS2021} take the value from the object-detection literature ``in practice''. They do compare focal loss against cross-entropy across a sweep of $\alpha$, but that experiment holds $\gamma$ fixed, so the focusing exponent remains unvalidated on this task.

\textbf{Label smoothing is named in the corpus but never examined on this task.} Dosovitskiy et al.~\cite{dosovitskiy2021} list it among the regularisation parameters they tune, but no reviewed micro-expression paper uses or discusses it, or reports what it does under class skew. It is used here and declared as a choice this corpus supports neither way.

\subsection{Implications for this research}
\label{sec:imbalance-double-correction}

\begin{table}[htbp]
	\centering
	\footnotesize
	\begin{tabular}{@{}p{5.6cm} p{8.0cm}@{}}
		\hline
		\textbf{Finding from the review} & \textbf{Decision taken in this thesis} \\
		\hline
		Focal loss addresses class-frequency bias in micro-expression recognition~\cite{zhaoS2021} & Focal loss adopted as the training objective, with \textbf{$\gamma = 2.0$}, matching Zhao et al.'s reported setting \\
		$\alpha$ is a dataset-specific class-weight term~\cite{zhaoS2021} & Implemented as an optional per-class $\alpha$ vector, \textbf{switched off at run time} --- see the declaration below \\
		Correct at the data end as well as the loss end~\cite{xia2020a} & Minority oversampling via a weighted sampler, active in every configuration --- a sampler, where Xia et al. use uniform augmentation plus a balanced loss \\
		Imbalance-aware metrics are mandatory~\cite{see2019,liong2016} & Pooled macro F1 --- identical in construction to MEGC's UF1 --- as the primary metric, with accuracy quoted only alongside the always-Negative floor \\
		\textit{(no micro-expression source)} & Label smoothing of 0.05 applied inside the focal loss; named among the regularisers tuned by Dosovitskiy et al.~\cite{dosovitskiy2021}, but declared as unsupported either way by this corpus \\
		\hline
	\end{tabular}
	\caption[Review findings and the imbalance-handling decisions they determine]{Review findings and the imbalance-handling decisions they determine.}
	\label{tab:imbalance-decisions}
\end{table}

\textbf{The finding this project contributes: correct at exactly one point, and know which.} The configuration is resolved at run time by a single rule --- class weighting in the loss is enabled only when the balanced sampler is \textit{not} active --- and the run log records the decision each fold: \textit{``Class weights in loss disabled (balanced sampler already active).''}

That rule exists because of two failures in opposite directions, both recoverable from this project's own run history. In the earliest code era the loss carried inverse-frequency class weights and no balanced sampler existed; under that single correction the proposed configuration abandoned the \textbf{Positive} class outright, predicting none of its clips --- a failure initially misdiagnosed as over-capacity overfitting. A later era added the sampler while leaving the loss weights unconditional, so both corrections ran at once; the collapse then inverted, and configurations abandoned the \textbf{Negative} majority class instead, one of them predicting a single class for every clip in the corpus. Correcting once at the loss starved the minority; correcting twice over-served it. With the guard in place and exactly one correction active, no configuration in the reported study abandons any class.

This is a contribution rather than an incident, because of how little guidance the corpus offers (\autoref{sec:imbalance-limitations}): no reviewed paper warns that combining corrections can invert the intended effect rather than merely blunt it. On a 156-clip corpus with a 4 : 1 skew, correcting twice is as damaging as correcting once in the wrong place --- the two failures differ only in which class is abandoned.

\textbf{A second declaration.} The severity of the imbalance also constrains what the ablation can conclude. Because every configuration shares the same sampler, loss and $\gamma$, this thesis measures \textbf{no} effect attributable to the imbalance treatment itself: it is a constant, not a factor. The claim supported is that the four architectural components were compared under a training regime that does not abandon minority classes, not that this regime is optimal.

\subsection{Research gap and scope}
\label{sec:imbalance-gap}

The corrections themselves are well established; what the corpus does not settle are the three gaps of \autoref{sec:imbalance-limitations} --- how two corrections compose, what $\gamma$ should be on this task, and what over-correction looks like when it fails.

This thesis contributes to the first and third, documenting two opposite class-collapse failures and the run-time rule that resolved them (\autoref{sec:imbalance-double-correction}): a negative result the corpus does not contain and one directly actionable for anyone reproducing this class of pipeline. On the second it offers no improvement, since $\gamma = 2.0$ and label smoothing of 0.05 are held constant across all twelve configurations. Sweeping either, and the Grad-CAM++ auditing and identity-disentanglement work this project's earlier stages removed, are declared as further work.
